% !TeX root = ../../Book3.tex
% 纲要标题：### H.1 曲线、除子与亏格
% 状态：新增正文初稿，尚未正式审核。

\section{曲线、除子与亏格}
\label{b3:sec:H01}

在仿射直线上，多项式只有零点；补入无穷远点以后，次数又成为那一点的极点阶数。例如坐标函数 \(t\) 在零点消失一次，在无穷远点有一阶极点。一般曲线未必有一个坐标足以表示全部函数，但每个光滑点仍有自己的统一化元。把这些局部阶数放在一起，就得到除子；规定函数所允许的极点，则得到一个线丛的截面空间。

\ProseParagraphBreak
本节在代数闭域 \(k\) 上工作，不限制特征。代数闭性使每个闭点的剩余域都等于 \(k\)，因而除子的次数可直接相加。若底域不代数闭，次数须按 \([\kappa(P):k]\) 加权，且一个次数一的除子类未必由有理点表示；后面关于亏格一曲线的点与除子类的对应依赖这里的基域条件。

\subsection{代数曲线与函数域}
\label{b3:subsec:H01:01}

\begin{definition}\label{b3:def:H-smooth-projective-curve}
	设 \(k\) 为代数闭域。一个一维、有限型、分离的整 \(k\) 概形称为本附录所讨论的\term[daishu-quxian]{代数曲线}[algebraic curve]。若它在 \(k\) 上光滑，并且存在闭嵌入 \(C\hookrightarrow\mathbb P_k^N\)，则称为光滑射影整曲线。其一般点记为 \(\eta\)，\(k(C)=\mathcal O_{C,\eta}\) 是它的\term[quxian-hanshuyu]{函数域}[function field]。
\end{definition}

任取非空仿射开集 \(U=\operatorname{Spec}A\)，便有 \(k(C)=\operatorname{Frac}(A)\)。这里 \(A\) 是有限型一维整环，故函数域在 \(k\) 上的超越次数为一。函数域中的元素是定义在某个非空开集上的正则函数，两个表示在共同定义域内相等就表示同一个有理函数。

\ProseParagraphBreak
代数闭域是完美域，\cref{b3:thm:perfect-field-regular-smooth}因而使光滑曲线在闭点 \(P\) 的局部环成为一维正则局部环。其极大理想由一个元素生成，\cref{b3:thm:dvr-characterizations}使 \(\mathcal O_{C,P}\) 成为 DVR。记相应归一化赋值为 \(\operatorname{ord}_P:k(C)^*\rightarrow\mathbb Z\)。它对函数给出零点或极点的阶数；统一化元的选择只改变单位因子。
\symidx{\(k(C)\)：整曲线 \(C\) 的函数域}
\symidx{\(\operatorname{ord}_P\)：光滑曲线闭点 \(P\) 处的阶数}

\begin{lemma}\label{b3:lem:H-rational-map-extension}
	设 \(k\) 为代数闭域，\(X\) 为正规整曲线，\(Y\subseteq\mathbb P_k^N\) 为射影闭子概形。每个从 \(X\) 到 \(Y\) 的有理 \(k\) 映射都唯一延拓为 \(k\) 态射 \(X\rightarrow Y\)。
\end{lemma}
\begin{proof}
	在一般点把映射写成齐次坐标 \([a_0:\cdots:a_N]\)，其中 \(a_i\in k(X)\)，且不全为零。给定闭点 \(P\)，其局部环为 DVR。取统一化元 \(t\)，令 \(m\) 为非零 \(a_i\) 的阶数最小值。将全部坐标乘以 \(t^{-m}\)，则每个坐标都在 \(\mathcal O_{X,P}\) 中，且至少一个为单位。缩小 \(P\) 的开邻域，使这些坐标都正则，选定的单位坐标处处不为零，就得到到射影空间的态射。\(Y\) 的齐次方程在函数域中成立，所以在这些邻域上仍成立，像落在 \(Y\)。

	两份局部延拓在一般点相等。因为 \(Y\) 分离，它们的等化子是闭子概形；其定义理想在整的源上泛点为零，因而为零。于是两态射相等，可唯一黏合。一般点原来就在定义域中，所有闭点又已覆盖，得到全体 \(X\) 上的延拓。
\end{proof}

这个坐标论证也见\cite[Chapter I, Proposition 6.8, pp. 43--44]{Hartshorne1977}。源必须有能把全部分母统一比较的局部参数；在奇异点，局部环未必是 DVR，向射影空间的有理映射就不能照搬这一消分母过程。

\begin{proposition}\label{b3:prop:H-curve-morphism}
	设 \(k\) 为代数闭域，\(C,D\) 为光滑射影整曲线，\(\varphi:C\rightarrow D\) 为非恒定 \(k\) 态射。则 \(\varphi\) 是有限、平坦、满射的态射。令
	\[
		d=[k(C):\varphi^*k(D)].
	\]
	对任一闭点 \(Q\in D\)，取其统一化元 \(u\)，则有
	\[
		\sum_{P\in\varphi^{-1}(Q)}\operatorname{ord}_P(\varphi^*u)=d.
	\]
	这个整数 \(d\) 称为态射的次数，包含函数域扩张的可分次数与不可分次数。若 \(d=1\)，则 \(\varphi\) 为同构。
\end{proposition}
\begin{proof}
	非恒定性使像稠密：一维整概形的真闭子集只有有限个闭点，整源的像若落在其中就只能是一个点。因此有函数域嵌入 \(k(D)\hookrightarrow k(C)\)。两域都在 \(k\) 上有限生成且超越次数一，所以此扩张代数且有限。

	取 \(D\) 的非空仿射开集 \(V=\operatorname{Spec}A\)，并令 \(B\) 为 \(A\) 在 \(k(C)\) 中的整闭包。\cref{b3:thm:normalization-finite}适用于有限型 \(k\) 整环 \(A\) 及这个有限域扩张，故 \(B\) 是有限 \(A\) 代数。有限域扩张中的元素清分母后，其某个非零 \(A\) 倍数整于 \(A\)，所以 \(\operatorname{Frac}(B)=k(C)\)。整扩张保持维数，故它是一维正规整环，\(\operatorname{Spec}B\) 的闭点局部环都是 DVR。一般点给出到 \(C\) 的有理映射，由上一引理延拓为 \(\psi:\operatorname{Spec}B\rightarrow C\)。复合 \(\varphi\psi\) 与自然映射 \(\operatorname{Spec}B\rightarrow D\) 在一般点相同，由分离性相同，故 \(\psi\) 的像落在 \(\varphi^{-1}(V)\)。

	反向，对 \(P\in\varphi^{-1}(V)\)，每个 \(b\in B\) 都整于 \(A\subseteq\mathcal O_{C,P}\)，故整于 \(\mathcal O_{C,P}\)。该局部环整闭，因此 \(b\in\mathcal O_{C,P}\)。这些有理函数在每一点正则，可黏合成 \(\varphi^{-1}(V)\) 上的正则函数，给出态射 \(\gamma:\varphi^{-1}(V)\rightarrow\operatorname{Spec}B\)。两复合都在一般点为恒等，再由分离性和源的整性，均为恒等。所以
	\[
		\varphi^{-1}(V)\cong\operatorname{Spec}B.
	\]
	这证明 \(\varphi\) 有限；整扩张的 lying-over 定理保证它在每张 \(V\) 上满射，故全局满射。

	在闭点 \(Q\) 处局部化，\(B_Q=B\otimes_A\mathcal O_{D,Q}\) 是 DVR \(\mathcal O_{D,Q}\) 上秩为 \(d\) 的有限生成无挠模，所以自由秩为 \(d\)。这一局部结论可直接用 Nakayama 引理检验：取模统一化元的一组基并提升为最小生成元；若存在非零关系，把系数中共同的最低幂约去，再模掉统一化元，就与这些剩余类线性无关矛盾。一般点处的局部映射是域扩张，也平坦。因此 \(\varphi\) 平坦，且
	\[
		\dim_k B_Q/uB_Q=d.
	\]
	这个有限维 \(k\) 代数分解为各闭点 \(P\mid Q\) 的局部因子的积。每个因子是 \(\mathcal O_{C,P}/(\varphi^*u)\)，其 \(k\) 维数由\cref{b3:prop:dvr-order-ideals}等于 \(\operatorname{ord}_P(\varphi^*u)\)。维数相加便得公式。最后，若 \(d=1\)，函数域相同；\(A\) 已整闭，所以 \(B=A\)，上述仿射识别说明 \(\varphi\) 是同构。
\end{proof}

有限态射与纤维重数公式也见\cite[Chapter II, Propositions 6.8--6.9, pp. 137--138]{Hartshorne1977}。例如特征 \(p>0\) 时，\(\mathbb P^1\rightarrow\mathbb P^1\)、\(t\mapsto t^p\) 的次数是 \(p\)。每个闭点纤维只有一个点，因为 \(t^p-a=(t-b)^p\)，其中 \(b^p=a\)；其重数却为 \(p\)。因此一般不能用不同纤维点的个数代替次数。

\subsection{除子与线丛}
\label{b3:subsec:H01:02}

\begin{definition}\label{b3:def:H-curve-divisor}
	设 \(k\) 为代数闭域，\(C\) 为光滑射影整曲线。一个除子是闭点的有限整数线性组合 \(D=\sum_P n_PP\)，次数为 \(\deg D=\sum_Pn_P\)。各系数非负时称为有效除子，记 \(D\ge0\)。对 \(f\in k(C)^*\)，定义
	\[
		\operatorname{div}(f)=\sum_P\operatorname{ord}_P(f)P.
	\]
	这样的除子称为主除子；若 \(D-E=\operatorname{div}(f)\)，则记 \(D\sim E\)，称 \(D,E\) 线性等价。
\end{definition}

主除子的支撑确实有限。取有限个仿射开集覆盖 \(C\)，在每张图上把 \(f\) 写成分式 \(a/b\)。非零的 \(a,b\) 在一维 Noether 整环中的零点只有有限个：包含它们的素理想都极小于相应主理想，而且都为闭点。有限张图的这些有限集合包含全部非零阶数点。赋值的加法公式使主除子组成除子群的子群，这与\secref{b3:sec:2003}的仿射定义相容。

\begin{corollary}\label{b3:cor:H-principal-degree}
	设 \(k\) 为代数闭域，\(C\) 为光滑射影整曲线。则每个主除子次数为零，并且 \(H^0(C,\mathcal O_C)=k\)。
\end{corollary}
\begin{proof}
	若 \(f\in k(C)^*\) 非恒定，它定义到 \(\mathbb P_k^1\) 的有理映射，由\cref{b3:lem:H-rational-map-extension}延拓为态射。零点纤维与无穷远纤维的总重数，由\cref{b3:prop:H-curve-morphism}都等于 \([k(C):k(f)]\)。前者组成 \(f\) 的零除子，后者组成极除子，故 \(\deg\operatorname{div}(f)=0\)。常函数的主除子为零，同样成立。

	全局正则函数没有极点。若它非恒定，相应映射 \(C\rightarrow\mathbb P^1\) 却必须满射，因而在无穷远点有非空纤维，与没有极点矛盾。因此全局函数都在 \(k\) 中；反向常函数当然全局正则。
\end{proof}

在 \(\mathbb P_k^1\) 上以 \(t\) 为仿射坐标，记无穷远点为 \(\infty\)。对不同的 \(a,b\in k\)，有
\[
	\operatorname{div}\!\left(\frac{(t-a)^2}{t-b}\right)=2[a]-[b]-[\infty].
\]
无穷远处令 \(s=t^{-1}\)，原函数为 \(s^{-1}(1-as)^2/(1-bs)\)，所以该处极点阶数为一。若只保留仿射直线，就漏掉这个负项；次数为零是射影曲线上的全局结论。

\begin{definition}\label{b3:def:H-divisor-sheaf}
	设 \(k\) 为代数闭域，\(C\) 为光滑射影整曲线，\(D=\sum n_PP\) 为除子。对非空开集 \(U\subseteq C\)，定义
	\[
		\mathcal O_C(D)(U)=\{0\}\cup\{\,f\in k(C)^*: \operatorname{ord}_P(f)+n_P\ge0\ \text{对所有闭点}\;P\in U\,\}.
	\]
	记 \(L(D)=H^0(C,\mathcal O_C(D))\)、\(\ell(D)=\dim_kL(D)\)。\(L(D)\) 是满足指定零点及极点界的函数空间。
\end{definition}
\symidx{\(\mathcal O_C(D)\)：除子 \(D\) 的关联线丛}
\symidx{\(L(D)\)、\(\ell(D)\)：除子截面空间及其维数}

这里采用几何的符号约定：正的 \(n_P\) 允许极点，而负的 \(n_P\) 要求零点。它与\secref{b3:sec:2005}中的分式理想 \(I_D\) 相差一个负号；那里的 \(I_D\) 对应这里的 \(\mathcal O_C(-D)\)。

\begin{proposition}\label{b3:prop:H-divisor-line-bundle}
	设 \(k\) 为代数闭域，\(C\) 为光滑射影整曲线，\(D,E\) 为 \(C\) 上的除子。则 \(\mathcal O_C(D)\) 为可逆层，且
	\[
		\mathcal O_C(D)\otimes\mathcal O_C(E)\cong\mathcal O_C(D+E).
	\]
	映射 \([D]\mapsto[\mathcal O_C(D)]\) 给出除子模线性等价所成的群与线丛同构类群 \(\operatorname{Pic}(C)\) 的同构。特别地，线丛的次数可定义为 \(\deg\mathcal O_C(D)=\deg D\)，与表示无关。
\end{proposition}
\begin{proof}
	在 \(P\) 处取统一化元 \(t\)，则所定义层的茎为 \(t^{-n_P}\mathcal O_{C,P}\)。缩小邻域，删去 \(D\) 的其他支撑点以及 \(t\) 的其他零点、极点，这个式子就给出邻域上的自由基。条件由局部正则性定义，因而满足限制与黏合。两个局部基相乘得到 \(t^{-n_P-m_P}\)，证明张量积公式。

	任取线丛 \(\mathcal L\)。其一般点纤维是一维 \(k(C)\) 空间，选非零有理截面 \(s\)。在有限组平凡化开集上写 \(s=f_ie_i\)，其中 \(e_i\) 是局部基。定义 \(n_P=\operatorname{ord}_P(f_i)\)；更换基只乘以单位，所以与选择的开集无关。各 \(f_i\) 的零极点有限，故 \(D_s=\sum n_PP\) 是除子。映射 \(h\mapsto hs\) 把 \(\mathcal O_C(D_s)\) 同构到 \(\mathcal L\)：在 \(P\) 处，\(t^{-n_P}s\) 是一个局部基。这证明满射性。

	若两线丛 \(\mathcal O_C(D),\mathcal O_C(E)\) 同构，其一般点上的映射是乘以某个 \(a\in k(C)^*\)。在 \(P\) 处比较两边的生成元阶数，得 \(\operatorname{ord}_P(a)-n_P(D)=-n_P(E)\)，即 \(D-E=\operatorname{div}(a)\)。反之，此等式使乘以 \(a\) 给出同构。因此核恰是主除子群。次数的良定义由\cref{b3:cor:H-principal-degree}得到。
\end{proof}

这也说明线丛上的非零正则截面给出有效零除子，且其次数等于线丛的次数；对 \(\mathcal O_C(D)\) 的非零截面 \(f\)，零除子正是 \(\operatorname{div}(f)+D\)。除子与可逆层的比较见\cite[Chapter II, Proposition 6.13, Proposition 6.15 and Corollary 6.16, pp. 144--145]{Hartshorne1977}。

\ProseParagraphBreak
次数是群同态 \(\deg:\operatorname{Pic}(C)\rightarrow\mathbb Z\)。其核记为 \(\operatorname{Pic}^0(C)\)，由零次数线丛的同构类组成，也可以写成零次数除子模主除子所成的群。这里把它作为抽象交换群使用；在亏格一时，\secref{b3:sec:H04}将把这个群结构转移到曲线的点集上。
\symidx{\(\operatorname{Pic}^0(C)\)：曲线 \(C\) 的零次数线丛类群}

\begin{corollary}\label{b3:cor:H-section-degree}
	设 \(k\) 为代数闭域，\(C\) 为光滑射影整曲线，\(\mathcal L\) 为线丛。若 \(\deg\mathcal L<0\)，则 \(H^0(C,\mathcal L)=0\)。若 \(\deg\mathcal L=0\)，则它有非零全局截面当且仅当 \(\mathcal L\cong\mathcal O_C\)，此时截面空间维数为一。
\end{corollary}
\begin{proof}
	非零截面的零除子有效，次数不可能为负。次数零时，这个有效除子只能为零，所以该截面处处生成线丛，给出 \(\mathcal O_C\cong\mathcal L\)。反向使用常数截面，再由\cref{b3:cor:H-principal-degree}计算维数。
\end{proof}

在射影直线上，\(\operatorname{div}(t-a)=[a]-[\infty]\)，所以每个除子都线性等价于 \((\deg D)[\infty]\)。不同次数不能线性等价，故 \(\operatorname{Pic}(\mathbb P^1)\cong\mathbb Z\)。在两张坐标图上，\(\mathcal O(n[\infty])\) 的基可取 \(1,t^n\)，其过渡关系正是\secref{b3:sec:2403}中 \(\mathcal O(n)\) 的过渡关系，因而二者同构。

\subsection{亏格与典范除子}
\label{b3:subsec:H01:03}

截面空间给出可以全局实现的函数，一次上同调则保留局部截面之间无法消去的差。为把它的维数作为曲线的数值，需要先确认维数有限。对曲线，这可以直接化为两个 Laurent 多项式格的比较。

\begin{proposition}\label{b3:prop:H-curve-cohomology-finite}
	设 \(k\) 为代数闭域，\(C\) 为光滑射影整曲线，\(\mathcal L\) 为线丛。则 \(H^0(C,\mathcal L)\)、\(H^1(C,\mathcal L)\) 都为有限维 \(k\) 空间，且 \(H^i(C,\mathcal L)=0\) 对 \(i\ge2\) 成立。
\end{proposition}
\begin{proof}
	选一个非恒定函数 \(t\in k(C)\)，它给出次数为 \(d\) 的有限态射 \(\varphi:C\rightarrow\mathbb P^1\)。两张标准图的逆像 \(U,V\) 都仿射，交集 \(W\) 也仿射。设 \(M=\Gamma(U,\mathcal L)\)、\(N=\Gamma(V,\mathcal L)\)。由\cref{b3:cor:vector-bundle-projective-module}，它们分别是两张仿射图坐标环上的有限生成投射模。因为 \(\varphi\) 有限且 \(\mathcal L\) 可逆，\(M,N\) 分别为 \(k[t]\)、\(k[t^{-1}]\) 上的有限生成无挠模；在函数域上，它们都给出一维 \(k(C)\) 空间，故在 \(k(t)\) 上的秩是 \(d\)。

% 跨卷引用目标：b2:thm:pid-module-structure（4.3.2）。
	第二卷定理4.3.2证明主理想整环上有限生成模的不变因子分解；这里两个多项式环都是 PID，无挠性排除全部循环挠商，因此 \(M,N\) 均自由秩为 \(d\)。在交集上限制截面对应于局部化，这是\cref{b3:thm:affine-quasicoherent-equivalence}的模层对应。选 \(M\) 的基，在 \(W\) 上得到 \(E=k[t,t^{-1}]^d\)。\(N\) 的基在其中由某个 \(T\in\operatorname{GL}_d(k[t,t^{-1}])\) 给出，所以
	\[
		M=k[t]^d,\qquad N=T\,k[t^{-1}]^d\quad\text{均作为}\;E\;\text{的子模}.
	\]
	由\cref{b3:thm:affine-cech-comparison}，上同调是两项复形 \(M\oplus N\rightarrow E\)、\((m,n)\mapsto n-m\) 的上同调。因此 \(H^0\cong M\cap N\)、\(H^1\cong E/(M+N)\)，更高次数为零。

	取非负整数 \(r\)，使 \(T\) 及 \(T^{-1}\) 的所有 Laurent 项指数均不超过 \(r\)。若 \(v\in M\cap N\)，写 \(v=Tb\)、\(b\in k[t^{-1}]^d\)，则 \(v\) 的最高指数不超过 \(r\)；又因 \(v\in M\)，其最低指数非负。于是 \(H^0\) 嵌入由 \(t^je_i\)、\(0\le j\le r\) 张成的有限维空间。对 \(H^1\)，所有 \(j\ge0\) 的 \(t^je_i\) 已属于 \(M\)；若 \(j\le-r\)，则 \(T^{-1}t^je_i\) 的指数都非正，故 \(t^je_i\in N\)。余核只需由 \(-r<j<0\) 的有限多个向量生成，也有限维。
\end{proof}

\begin{definition}\label{b3:def:H-genus-canonical}
	设 \(k\) 为代数闭域，\(C\) 为光滑射影整曲线。定义它的\term[kuige]{亏格}[genus]为
	\[
		g(C)=\dim_kH^1(C,\mathcal O_C).
	\]
	将各仿射图上的微分模黏合，得到微分层 \(\Omega_{C/k}\)，记为 \(\omega_C\)，称为曲线的\term[dianfan-xiancong]{典范线丛}[canonical line bundle]。任取非零有理微分 \(\alpha\in\Omega_{k(C)/k}\)，其作为 \(\omega_C\) 的有理截面的除子称为\term[dianfan-chuzi]{典范除子}[canonical divisor]，记为 \(K_C\)。
\end{definition}
\symidx{\(g(C)\)：光滑射影整曲线的亏格}
\symidx{\(\omega_C\)、\(K_C\)：典范线丛与典范除子}

微分层的黏合由\cref{b3:prop:differential-localization}保证；\cref{b3:prop:standard-smooth-lifting,b3:thm:smooth-local-standard-form}使其在光滑曲线上局部自由秩一。在闭点处，统一化元 \(t\) 的类生成 \(\mathfrak m_P/\mathfrak m_P^2\)，由\cref{b3:prop:tangent-cotangent-rational}，\(\mathrm{d}t\) 模极大理想非零，故生成 \(\omega_{C,P}\)。写 \(\alpha=a\,\mathrm{d}t\)，便有 \(\operatorname{ord}_P(\alpha)=\operatorname{ord}_P(a)\)。更换统一化元时，两组局部基相差单位，阶数不变。更换非零有理微分 \(\alpha\) 为 \(f\alpha\)，典范除子增加 \(\operatorname{div}(f)\)；因此 \(K_C\) 只在具体代表上有选择，线性等价类固定，且 \(\omega_C\cong\mathcal O_C(K_C)\)。

\ProseParagraphBreak
在正特征中不能把任意非恒定函数的微分当作典范除子的代表。例如 \(\mathrm{d}(t^p)=0\)，零微分没有定义有限阶数的除子。上面的定义明确要求非零有理微分；光滑性保证函数域上的微分空间是一维，因而存在这样的选择。

\begin{example}\label{b3:exm:H-projective-line-canonical}
	设 \(k\) 为代数闭域，\(C=\mathbb P_k^1\)。在有限点 \(a\)，\(t-a\) 为统一化元，且 \(\mathrm{d}t=\mathrm{d}(t-a)\)，所以 \(\mathrm{d}t\) 没有零极点。在无穷远处令 \(s=t^{-1}\)，则
	\[
		\mathrm{d}t=-s^{-2}\,\mathrm{d}s.
	\]
	所以 \(\operatorname{div}(\mathrm{d}t)=-2[\infty]\)，并有 \(\omega_{\mathbb P^1}\cong\mathcal O(-2)\)。这个等式在特征二仍成立，因为系数 \(-1\) 仍为单位。由\cref{b3:prop:projective-line-cohomology}，\(g(\mathbb P^1)=h^1(\mathcal O)=0\)。同时 \(h^0(\omega_{\mathbb P^1})=0\)，这是\secref{b3:sec:H03}中一般对偶关系的一个具体情形。
\end{example}

亏格目前定义为 \(h^1(\mathcal O_C)\)，而 \(h^0(\omega_C)\) 数的是正则微分。它们相等并不是定义的一部分；\cref{b3:thm:H-serre-duality}将通过一个非退化配对建立两者的联系。典范除子的次数也由\secref{b3:sec:H02}的截面维数公式确定。

% BEGIN APPENDIX OUTLINE SYNC
% * 状态：初稿已完成；新增范围尚未正式审核。
% #### H.1.1 代数曲线与函数域
% #### H.1.2 除子与线丛
% #### H.1.3 亏格与典范除子
% END APPENDIX OUTLINE SYNC
